\section{Research Context}
% =============================================================================
% RESEARCH CONTEXT - v3 EMERGENT SOCIAL STRUCTURES FRAMING
% =============================================================================
% STATUS: Needs restructure around cooperation literature, not IR theory
%
% FRAMING EVOLUTION:
%   v1: Cold War → Waltz → Wendt → both valid (IR-heavy)
%   v2: Security dilemma → two hypotheses → gap (narrow)
%   v3: Cooperation literature → semantic priors → emergent structures (rich)
%
% NEW STRUCTURE:
%   Para 1: COOPERATION IS WELL-STUDIED
%           - Axelrod (1984): optimizers can achieve cooperation
%           - Game theory, ABM tradition
%           - BUT: this cooperation is mechanistic, fragile, slow
%
%   Para 2: SEMANTIC REASONING IS DIFFERENT
%           - LLMs bring something new: semantic priors
%           - Norm recognition, identity formation, narrative interpretation
%           - QUESTION: Do these enable qualitatively different structures?
%
%   Para 3: THE GAP
%           - RL vs LLM comparisons exist but don't isolate variables
%           - No systematic study of whether semantic priors produce
%             different EMERGENT STRUCTURES (not just different outcomes)
%           - Your contribution: controlled comparison, same goal, different architecture
%
%   Para 4: RESEARCH QUESTION
%           - "Do semantic priors enable qualitatively different emergent
%              social structures than pure optimization?"
%           - Measure: coalitions, cooperation, hierarchy, conflict, norms
%
% OPTIONAL: Brief IR context as APPLICATION DOMAIN (1-2 sentences max)
%   "This question has implications for understanding cooperation in
%   international relations, where the neorealism-constructivism debate
%   mirrors the optimization vs semantic reasoning tension."
%
% CITATIONS TO USE:
%   - \cite{axelrod1984} - cooperation among self-interested agents
%   - \cite{cederman1997} - ABM tradition
%   - \cite{hua2023waragent} - LLM agents in strategic settings
%   - \cite{pan2025} - RL agents
%   - \cite{larooij2025} - generative ABM methodology
%
% EXAMPLE DIRECTION:
%
%   Para 1: "Cooperation among self-interested agents is a central problem
%   in multi-agent systems. Axelrod (1984) demonstrated that optimizing agents
%   can achieve cooperation through repeated interaction, discovering strategies
%   like tit-for-tat. This work spawned extensive research in game theory and
%   agent-based modeling. However, the cooperation that emerges is fundamentally
%   mechanistic: agents stumble into stable equilibria without 'understanding'
%   cooperation."
%
%   Para 2: "Large Language Models introduce a qualitatively different decision
%   process. Unlike optimization-based agents, LLMs bring semantic priors: the
%   ability to recognize norms, form identities, interpret actions through
%   narratives. When an LLM agent sees another agent invest in it, it doesn't
%   just update a value function - it can interpret the action as trust,
%   friendship, or alliance-building."
%
%   Para 3: "Recent work explores both RL agents (Pan 2025) and LLM agents
%   (Hua 2023) in strategic settings. However, these studies typically compare
%   agents with different goals, conflating architecture and objective. No study
%   systematically tests whether semantic reasoning produces different emergent
%   social structures when agents share the same optimization goal."
%
%   Para 4: "This thesis asks: Do semantic priors enable qualitatively different
%   emergent social structures than pure optimization? We compare self-optimizing
%   RL agents with self-optimizing LLM agents, measuring coalition formation,
%   cooperation patterns, hierarchy emergence, and conflict dynamics."
% =============================================================================

The ending of the Cold War introduced an intellectual vacuum as existing paradigms were unable to predict its peaceful conclusion. \textbf{Neorealism}, as outlined by Kenneth Waltz \cite{waltz1979}, theorizes that states act out of self-interest, be it for survival or power accumulation, given the anarchical system that is the world stage. The USSR undermining their relative gain to the US could be explained through the lens of upcoming \textbf{Constructivism}, formalized by Alexander Wendt \cite{wendt1992} in 1992 as a direct criticism on Neorealism. Instead of being rigidly self-interested through material constraints, state identities and interests are formed through interaction and are changeable. As of today, both theories are recognized to have explanatory abilities.

The utilization of ABM within the IR field has extensive applications \cite{cederman1997}. More recently both the use of RL agents \cite{pan2025} and LLM agents \cite{hua2023waragent} has been explored. However, no direct comparison between Neorealism and Constructivism has been done and computational mapping both framework to respectively RL agents and LLM agents is novel.

Research Question:
\textit{Under what material constraints do semantic reasoning (LLM agents) and
pure optimization (RL agents) produce divergent emergent social structures in
an anarchic system?}

% TODO: Complete rewrite around cooperation literature:
%
% STRUCTURE:
%   1. Cooperation is well-studied (Axelrod) but mechanistic
%   2. LLMs bring semantic priors (norms, identity, narrative)
%   3. Gap: no controlled comparison of emergent structures
%   4. RQ: Do semantic priors enable different emergent structures?
%
% KEEP IR BRIEF: 1-2 sentences as application domain, not main framing
%
% NEW RQ PHRASING:
%   "Do semantic priors enable qualitatively different emergent social
%   structures than pure optimization, or do material incentives dominate
%   regardless of reasoning architecture?"
